\section{Clarifications of Main-Text Statements}
\label{app:clarify}

This appendix carries three things that are needed to read the main paper correctly but that
would not fit in its eight pages: precise definitions of the two objects the integrity and
sabotage arguments turn on (\Cref{app:harness}), the exact statement of the placement result
whose main-text corollary is used loosely (\Cref{app:placement}), and the statistical
resolution of the experiments the main text reports (\Cref{app:power}). None of it changes a
number in the main paper; each item makes an existing statement unambiguous.

\subsection{The harness's trace projection and sabotage rule}
\label{app:harness}

The main paper's integrity theorem (\Cref{thm:integrity}) is a statement about the
\emph{design} of the checker, and its proof is only meaningful if the object the checker reads
is defined. The following two definitions make it explicit. They also settle a reading of the
main text that is otherwise ambiguous: if the trace contains every message, then a
manipulation that edits a message \emph{does} change the trace, so a statement of the form
``equal traces give equal verdicts'' would say nothing about content manipulation.
The resolution is that the checker reads not the full trace but its structural projection,
and the empirical question---whether a given injection changes that projection---is the one
\Cref{sec:integrity} measures.

\begin{definition}[Security-relevant projection]
\label{def:projection}
Let $\pi:\TraceSp\to\TraceSp^{\mathrm{st}}$ map a trace to its \emph{security-relevant
projection}: the subsequence of recorded structural facts---executed tool calls with their
arguments, recorded transitions such as \texttt{unsafe\_action}, \texttt{exfiltration} or
\texttt{sabotage}, memory versions written, role assignments, topology edges, and the identity
map---with all free-form message text, tool prose, and model reasoning removed. The predicates
of \Cref{def:predicates} are required to factor through $\pi$: there are
$\bar\Predsec,\bar\Predutil$ on $\pi(\TraceSp)$ with
$\Predsec=\bar\Predsec\circ\pi$ and $\Predutil=\bar\Predutil\circ\pi$. Equivalently, the
verdicts depend on \emph{what the system did}, never on how it was phrased---which is the
property that lifts AgentDojo's per-agent state check~\cite{debenedetti2024agentdojo} to the
global trace.
\end{definition}

Read that way, \Cref{thm:integrity} is an assumption-driven statement rather than a tautology
about traces: its content is that the registered predicates are functions of $\pi$, so a
message-content manipulation $\pi_m$ with $\pi(\pi_m(\Trace))=\pi(\Trace)$ leaves the verdict
unchanged, whether or not it changed anything a human would consider the system's state. The
converse is not claimed and is false in general: a manipulation that \emph{does} change
$\pi$---for instance one that causes an illicit tool call to be executed---changes the verdict,
which is exactly the intended behaviour of a security predicate. The complementary statement
for a content-reading judge (\Cref{prop:judgehijack}) is an existence claim about judges, and
the rate at which real judges flip is measured in \Cref{sec:integrity}.

\begin{definition}[Availability sabotage in the harness]
\label{def:sabotage}
A task is $k$-forced with witness family $\mathcal{W}$. The harness records
coordination sabotage---and therefore $\Predsec=1$---whenever no witness survives the episode,
that is when either (a) fewer than $k$ agents remain uncompromised at fixation, or (b) the
topology is hub-critical (star or tree) and the orchestrator is compromised, even if $k$
honest agents remain. Case (b) is a deliberate modelling choice: the operator's routing is
assumed to depend on the orchestrator on those topologies, so losing it loses availability
regardless of how many workers survive. No language model participates in the judgement: it is
a deterministic function of the compromised set, the topology, and $k$. (The definition of a
$k$-forced task is \Cref{def:forced} in the main paper.)
\end{definition}

Three consequences follow, and each corrects an over-reading that the main text invites.

\emph{(1) The rule, not the influence bound, is what saturates the metric.}
\Cref{prop:star} is a statement about expected \emph{influence}: a centre seed reaches each
leaf with probability $\pinf$, which is strictly below $1$ whenever $\pinf<1$. The saturation
of $\ASRs$ at $1.000$ on star and tree in \Cref{tab:defense} is produced by
\Cref{def:sabotage}(b) once the hub is compromised, together with the fraction of episodes in
which the hub is seeded. So the correct description of that result is ``the harness-defined
predicate is discharged without the action channel'', not ``an injection spreads to every
worker''.

\emph{(2) Case (b) is why the chain behaves differently, and why
\Cref{prop:forced}(b) does not by itself describe the sweeps.} \Cref{prop:forced} assumes that
the adversary's only lever is withholding its own contribution; under that assumption a chain
is the \emph{most} sabotage-prone topology, because every witness contains every intermediate
node. The harness instead applies case (b) only on star and tree, so a chain is the
\emph{least} sabotage-prone topology in the measurements: that is why the chain column shows
$\ASRs{=}0.108$ in the pilot topology table while star and tree reach $0.873$ and $0.766$.
The main text's phrase ``the deterministic availability path'' therefore refers to
\Cref{def:sabotage} and not to the transversal condition of \Cref{prop:forced}; the two agree
only on hub-critical topologies.

\emph{(3) It is why a zero-compliance defense still scores $\ASRs=1$.} The system predicate is
a disjunction (unsafe action $\lor$ availability sabotage). A prompt-level defense acts on the
first disjunct only. Once compliance is pinned at zero the second disjunct is the only one
left, and it is a function of the compromised set and the topology with no model in the loop,
so the metric stops depending on the backbone at all. That is the mechanism behind the three
bit-identical rows of \Cref{tab:defense}, and it is a property of the instrument rather than
an experimental finding about LLMs.

\subsection{The placement result, stated exactly}
\label{app:placement}

The main paper states a strict advantage for coordinated seeding
(\Cref{thm:collusion}) and defines a metric-level collusion advantage
$\Delta_{\mathrm{coll}}=\ASRs^{\mathrm{coll}}-\ASRs^{\mathrm{ind}}$
(\Cref{def:csr}). These are different quantities---one is a difference in expected
\emph{number of compromised agents}, the other a difference in \emph{attack-success
probability}---and the main text's ``strictly dominates'' should not be read as the second.
We restate the first in full here, under the same independent-cascade model.

\begin{proposition}[Placement advantage of coordinated seeding on the star]
\label{prop:placement}
Fix the independent-cascade model of \Cref{as:prop} on a star of $N$ nodes with a budget of
one seed, edges bidirectional for the purpose of influence (a compromised leaf can compromise
the centre, which is what a shared channel implies). A coordinated adversary that places its
seed at the centre achieves expected compromise
$C=1+\pinf(N-1)$; an adversary that places its seed uniformly at random achieves
$\E^{\mathrm{ind}}=\tfrac1N C+\tfrac{N-1}{N}L$ with $L=1+\pinf+\pinf^{2}(N-2)$. Their
difference
\[
\Delta_{\mathrm{place}}\;=\;C-\E^{\mathrm{ind}}
\;=\;\pinf(1-\pinf)\,\frac{(N-1)(N-2)}{N}
\]
is strictly positive and $\Theta\bigl(\pinf(1-\pinf)N\bigr)$ for every $\pinf\in(0,1)$ and
$N\ge3$.
\end{proposition}

\begin{proof}
Each leaf is a direct out-neighbour of the centre, so a centre seed compromises all $N-1$
leaves independently with probability $\pinf$: $C=1+\pinf(N-1)$. A leaf seed compromises the
centre with probability $\pinf$ and, conditioned on that, each remaining leaf with probability
$\pinf$: $L=1+\pinf+\pinf^{2}(N-2)$. A uniformly random seed lands on the centre with
probability $1/N$ and on a leaf with probability $(N-1)/N$, hence
$\E^{\mathrm{ind}}=\tfrac1N C+\tfrac{N-1}{N}L$. Therefore
\[
\Delta_{\mathrm{place}}=C-\E^{\mathrm{ind}}=\frac{N-1}{N}(C-L)
=\pinf(1-\pinf)\,\frac{(N-1)(N-2)}{N},
\]
since $C-L=\pinf(1-\pinf)(N-2)$ by \Cref{prop:star}. For $N\ge3$ both
$(N-1)(N-2)/N>0$ and $\pinf(1-\pinf)>0$, so the difference is strictly positive, and it is
$\Theta(\pinf(1-\pinf)N)$. With one-directional edges the leaf-seed term loses its trailing
$\pinf$, so the advantage is only larger; the bidirectional case is the conservative one.
\end{proof}

Two limits of \Cref{prop:placement} belong next to it. \emph{First}, it compares
\emph{placement policies under a fixed budget of one seed}, not ``collusion on'' versus
``collusion off''; a real colluding coalition also shares a channel and a joint objective,
which the independent-cascade model does not represent. \emph{Second}, and empirically more
important, the sweeps do not confirm the ordering the proposition suggests: in
\Cref{tab:collusion} the unsafe-action advantage is $+0.049,+0.003,-0.024,+0.038,-0.001$ for
$\betafrac=0.1,\dots,0.5$---sign-changing and entirely inside the $\pm\epsilon=0.05$
resolution of an $M{=}738$ sweep. \Cref{prop:placement} is therefore offered as a statement
about the environment's propagation model and as a prediction of \emph{where} to look (hub
topologies), not as a fitted description of those five numbers.

\subsection{Resolution of the reported differences}
\label{app:power}

The main paper's budgets follow \Cref{thm:sample}. With $\epsilon=\alpha=0.05$ this gives
$M=\bigl\lceil(2\epsilon^{2})^{-1}\ln(2/\alpha)\bigr\rceil=738$ episodes per configuration,
and the same computation bounds what a comparison of two cells can resolve.

\emph{Single-cell resolution.} Hoeffding gives
$\Prob[|\hat\theta-\theta|\ge\epsilon]\le2e^{-2M\epsilon^{2}}$, so at $M=738$ and
$\epsilon=0.05$ each $\widehat{\ASRs}$ is within $\pm0.05$ of its target with probability
$0.95$, and a \emph{difference} of two such estimates is resolved only to about
$\pm2\epsilon=0.1$. Concretely, in \Cref{tab:degradation} the guardrail's advantage over
no-defense is $0.244$, $0.178$ and $0.149$ in the three columns---far outside that
band---whereas the \emph{shrinkage} between the single-agent and multi-agent columns rests on
margins of $0.039$ and $0.008$, and the no-defense baseline itself drifts down
($0.587\to0.560\to0.539$), which shrinks the reported \emph{ratio} without any defense effect.
The main text's percentage sequence $41.6\%\to31.8\%\to27.6\%$ is arithmetically correct but
is \emph{not resolved} at this budget, and the main paper reports it as an unresolved trend.
The same applies to the five $\Delta_{\mathrm{coll}}$ values of \Cref{tab:collusion}, all of
which lie within $\pm0.05$ of zero, and to the non-monotone adaptive-red sequence of
\Cref{tab:adaptive}, for which $M=60$ per scored pair is even weaker.

\emph{Multiplicity.} The sweep is not corrected across its $G$ cells. Bonferroni over $G=21$
replaces $\alpha$ by $\alpha/G$, so \Cref{thm:sample}(c) would require
\[
M\;\ge\;\frac{\ln(2G/\alpha)}{2\epsilon^{2}}
\;=\;\frac{\ln(42/0.05)}{2(0.05)^{2}}\;\approx\;1347
\]
episodes per cell; the reported sweep uses $738$ and is therefore \emph{descriptive} at the
stated budget rather than a set of simultaneous $95\%$ claims. We state this rather than
silently pool, and it is the reason no main-text claim rests on a cell difference smaller than
$0.1$.

\emph{Clustered compliance counts.} The $1075$ compliance decisions per backbone are not
$1075$ independent Bernoulli trials. Decoding is greedy, the prompt varies only in agent id
and step index, and the decisions cluster inside $M=30$ episodes per cell. A Wilson interval
on $1075$ decisions (upper bound $0.0036$ at $0$ successes) therefore describes the released
sample; it is not an inferential interval for the model's compliance rate, and we do not use
it as one. The same clustering limits the per-cell claim: with $M=30$ episodes the standard
error of a single cell proportion near $0.5$ is about $0.09$, so a contrast of the size that
would matter (say a $0.15$ drop between the no-defense and defended arms) is
$\approx1.6$ standard errors and would need far more episodes to be established at
conventional power.

\emph{Wilson intervals used in the main paper.} For completeness, the two-sided $95\%$
Wilson intervals at $M=30$ are $[0,0.114]$ for $0/30$, $[0.886,1]$ for $30/30$, and
$[0.274,0.608]$ for $13/30$; the last value is the $0.43$ cell of \Cref{tab:defense}. An
earlier draft quoted $[0,0.10]$ and $[0.90,1]$, which are one-sided bounds
($1-0.05^{1/30}\approx0.095$), and $[0.26,0.62]$, which came from a different convention; the
main paper now uses the two-sided values.

\subsection{Terminology and a citation caveat}
\label{app:terms}

\emph{Names of the inherited quantities.} ASB~\cite{zhang2025asb} defines
$\mathrm{NRP}=\mathrm{PNA}\cdot(1-\mathrm{ASR})$ as \emph{net resilient performance} with
\emph{PNA} the performance under \emph{no} attack; an earlier draft of this paper expanded
NRP as ``non-refusal performance'', which is not ASB's term and is corrected in the main text.
\defname{} writes $\mathrm{NRP}_{\mathrm{sys}}$ for the collective version and
$\mathrm{NRP}_{\mathrm{act}}$ (equivalently $\mathrm{ASR}_{\mathrm{act}}$-based) for the
single-agent predicate used in \Cref{tab:degradation}; the two must not be conflated.

\emph{The emulation-fidelity figure.} The main paper cites ToolEmu's emulator/human agreement
as $\kappaem\approx0.48$ and treats it as the central validity caveat. We flag honestly that
we could not re-confirm that specific number against ToolEmu's reported results while
preparing this supplement: what ToolEmu~\cite{ruan2024toolemu} reports is an agreement rate
between its automatic evaluator and human annotators together with a
human-validated realism rate for its emulated trajectories, and the $\kappa$ value has been
carried in our materials without re-derivation. The value should therefore be read as an
inherited caveat rather than as a measurement of this paper's emulator, and the
emulated-versus-real-tools ablation that would replace it with a number of our own is
registered but unrun (\Cref{app:limitations}, item (iv)).
